\BourbakiUnnumberedChapter{Bibliography}{Bibliography}

\begin{enumerate}
\def\labelenumi{(\Roman{enumi})}
\tightlist
\item
  Euclidis Elementa, 5 vol., ed.~J. L. Heiberg, Lipsiae (Teubner), 1883-88.
\end{enumerate}

(I bis) T. L. HEATH, The thirteen books of Euclid's Elements\ldots, 3 vol., Cambridge, 1908.

\begin{enumerate}
\def\labelenumi{(\Roman{enumi})}
\setcounter{enumi}{1}
\tightlist
\item
  Diophanti Alexandrini Opera Omnia\ldots, 2 vol., ed.~P. Tannery, Lipsiae (Teubner), 1893-95.
\end{enumerate}

(II bis) T. L. HEATH, Diophantus of Alexandria, 2 \(^{nd}\) ed., Cambridge, 1910.

\begin{enumerate}
\def\labelenumi{(\Roman{enumi})}
\setcounter{enumi}{2}
\item
  B. DATTA and A. N. SINGH, History of Hindu Mathematics, 2 vol., Lahore (Motilal Banarsi Das), 1935-38.
\item
  S. STEVIN, Les œuvres mathématiques\ldots, ed.~A. Girard, Leyde (Elsevier), 1634, vol.~I.
\item
  L. EULER: a) Introductio in Analysin Infinitorum (Opera Omnia, (1), t. IX, Zürich-Leipzig-Berlin (O. Füssli and B. G. Teubner), 1945, p.~384); b) Theoria motus corporum solidorum seu rigidorum (Opera Omnia (2), t. III, Zürich-Leipzig-Berlin (O. Füssli and B. G. Teubner), 1948, p.~200-201); c) Vollständige Anleitung zur Algebra (Opera Omnia (1), t. I, Leipzig-Berlin (Teubner), 1911, p.~422); d) Recherches sur les racines imaginaires des équations (Opera Omnia (1), t. VI, Leipzig-Berlin (Teubner), 1921, p.~78).
\item
  J.-L. LAGRANGE, Œuvres, Paris (Gauthier-Villars), 1867-1892: a) Solutions de divers problèmes de Calcul intégral, t. I, p.~520; b) Recherches sur les équations séculaires du mouvement des nœuds, t. VI, p.~655-666; c) Recherches d'arithmétique, t. III, p.~695-795; d) Sur la forme des racines imaginaires des équations, t. III, p.~479; e) Nouvelle solution du problème de rotation d'un corps quelconque qui n'est animé par aucune force accélératrice, t. III, p.~579-616.
\item
  P. S. LAPLACE : a) Mémoire sur les solutions particulières des équations différentielles et sur les inégalités séculaires des planètes (Œuvres, t. VIII, Paris (Gauthier-Villars), 1891, p.~325-366) ; b) Mémoire sur les inégalités séculaires des planètes et des satellites (Œuvres, t. XI, Paris (Gauthier-Villars), 1895, p.~49-92).
\item
  C. F. GAUSS, Werke, t. I (Göttingen, 1870), t. II (ibid., 1876) et t. III (ibid., 1876).
\end{enumerate}

(VIII bis) Die vier Gauss'schen Beweise für die Zerlegung ganzer algebraischer Functionen in reelle Factoren ersten oder zweiten Grades (Ostwald's Klassiker, n° 14, Leipzig (Teubner), 1904).

\begin{enumerate}
\def\labelenumi{(\Roman{enumi})}
\setcounter{enumi}{8}
\item
  N. H. ABEL, Œuvres, t. I, ed.~Sylow and Lie, Christiania, 1881.
\item
  A. L. CAUCHY : a) Leçons sur les applications du Calcul infinitésimal à la Géométrie (Œuvres complètes (2), t. V, Paris (Gauthier-Villars), 1903, p.~248) ; b) Sur l'équation à l'aide de laquelle on détermine les inégalités séculaires des planètes (Œuvres complètes (2), t. IX, Paris (Gauthier-Villars), 1891, p.~174).
\item
  P. G. LEJEUNE-DIRICHLET, Werke, t. I, Berlin (G. Reimer), 1889, p.~619-644.
\item
  E. KUMMER, Zur Theorie der complexen Zahlen, J. de Crelle, t. XLIII (1847), p.~319 (Collected papers, vol.~I, Heidelberg (Springer V.), 1975, p.~203).
\item
  Ch. HERMITE, Œuvres, t. I, Paris (Gauthier-Villars), 1905.
\item
  J. J. SYLVESTER, Collected Mathematical Papers, vol.~I, Cambridge, 1904: An enumeration of the contacts of lines and surfaces of the second order, p.~219 (= Phil. Mag., 1851).
\item
  W. R. HAMILTON, Lectures on Quaternions, Dublin, 1853.
\item
  A. CAYLEY, Collected Mathematical Papers, Cambridge, 1889-1898: A memoir on the theory of matrices, t. II, p.~475-496 (= Phil. Trans., 1858).
\item
  H. J. SMITH, Collected Mathematical Papers, vol.~I, Oxford, 1894; On systems of linear indeterminate equations and congruences, p.~367 (Phil. Trans., 1861).
\item
  E. SCHERING, Die fundamental Classen der zusammensetzbaren arithmetischen Formen, Abh. Ges. Göttingen, t. XIV (1868-69), p.~13.
\item
  K. WEIERSTRASS, Mathematische Werke, Bd. II, Berlin (Mayer und Müller), 1895: Zur Theorie der bilinearen und quadratischen Formen, p.~19.
\item
  L. KRONECKER, Auseinandersetzungen einiger Eigenschaften der Klassenanzahl idealer complexer Zahlen, Monats. Abhandl. Berlin (1870), p.~881 (= Werke, t. I, Leipzig (Teubner), 1895, p.~273).
\item
  C. JORDAN, Traité des substitutions et des équations algébriques. Paris (Gauthier-Villars), 1870, p.~114-125.
\item
  G. FROBENIUS, Theorie der linearen Formen mit ganzen Coefficienten, Gesammelte Abhandlungen, vol.~I, Heidelberg (Springer V.), 1968, p.~482 (= J. de Crelle, 1879).
\item
  G. FROBENIUS und L. STICKELBERGER, Ueber Gruppen von vertauschbaren Elementen, J. de Crelle, t. LXXXVI (1879), p.~217 (= Frobenius, Ges. Abh., vol.~I, p.~545).
\item
  R. DEDEKIND, Gesammelte mathematische Werke, t. II, Braunschweig (Vieweg), 1932: Ueber Zerlegungen von Zahlen durch ihre grössten gemeinsamen Teiler, p.~103.
\item
  \begin{enumerate}
  \def\labelenumii{\alph{enumii})}
  \tightlist
  \item
    E. ARTIN und O. SCHREIER, Algebraische Konstruktion reeler Körper, Abh. Math. Sem. Univ. Hamburg, t. V (1927), p.~83; b) E. ARTIN, Ueber die Zerlegung definite Funktionen in Quadrate (ibid., p.~100); c) E. ARTIN und O. SCHREIER, Eine Kennzeichnung der reell abgeschlossenen Körper (ibid., p.~225).
  \end{enumerate}
\end{enumerate}

